\scanpage{34}
% archive sheet number: 33
\begingroup
\makeatletter\tagsleft@true\makeatother

Posons
\[
\Sigma_{12}=\Sigma_1^-\cap\Sigma_2^-,\qquad
E_1=\Sigma_1^+\cap\Sigma_1^-,\qquad
E_2=\Sigma_1^+\cap\Sigma_2^-\tag*{(13)}
\]
alors on a le diagramme d'inclusion de parties fermées de $X$

\usetikzlibrary{decorations.pathmorphing}
\diagram{\[
\begin{tikzcd}[ampersand replacement=\&,row sep=1.1em,column sep=.4em]
\& \& X \& \& \\
\& \Sigma_2^- \arrow[ur] \& \& \Sigma_1^- \arrow[ul] \& \\
\Sigma_1^+ \arrow[ur] \& \& \Sigma_{12} \arrow[ul]\arrow[ur] \& \& \Sigma_2^+ \arrow[ul]\\
\& E_1 \arrow[ul]\arrow[ur] \& \& E_2 \arrow[ul]\arrow[ur] \& \\
\& \& \unsure{\deleted{\Omega|2|\varnothing}}
\arrow[ul,shorten >=.8em]
\arrow[ul,no head,shorten >=.8em,decorate,decoration={zigzag,segment length=3.4pt,amplitude=1.6pt}]
\& \&
\end{tikzcd}\tag*{(14)}
\]}

\margin{NB $\Sigma_1^+\cap\Sigma_2^+=\varnothing$ \unsure{donc|donne}
$E_1\cap\Sigma_2^+=\varnothing$, $E_2\cap\Sigma_1^-=\varnothing$
i.e. $E_1\subset\Sigma_2^{\unsure{\circ|o|0}-}$, $E_2\subset\Sigma_1^{\circ-}$
et $E_1\cap E_2\unsure{.|-}=\varnothing$}

où les trois carrés sont cartésiens et cocartésiens
\[
\deleted{\unsure{E_1|E_+}}
\]
i.e.
\[
\left|\begin{aligned}
&E_1=\Sigma_1^+\cap\Sigma_{12},\quad \Sigma_2^-=\Sigma_1^+\cup\Sigma_{12}\\
&\deleted{E_1}
\end{aligned}\right.
\]
et les deux relations analogues correspondant aux deux autres carrés
\[
\left\{\begin{aligned}
&E_2=\Sigma_{\deleted{\unsure{2|1}}12}^{\deleted{\unsure{+|-}}}\cap\Sigma_{\deleted{\unsure{12|1}}2}^+,\quad
\Sigma_1^-=\Sigma_{22}^{\deleted{\unsure{+|-}}}\cup\Sigma_2^+\\
&\Sigma_{12}\overset{(13)}{\equiv}\Sigma_1^-\cap\Sigma_{\unsure{2|1}}^-,\quad
\unsure{\deleted{-|\Sigma}}X=\Sigma_2^-\cup\Sigma_1^-.
\end{aligned}\right.\tag*{(15)}
\]

Notons que
\[
\left\{\begin{aligned}
&\text{a)}\quad\Sigma_1^+\cap\Sigma_2^-=\Sigma_1^+\ne\varnothing\\
&\text{b)}\quad\Sigma_2^+\cap\Sigma_1^-=\Sigma_2^+\ne\varnothing\\
&\text{c)}\quad\Sigma_1^-\cap\Sigma_2^-=\Sigma_{12}\ne\varnothing,
\end{aligned}\right.\tag*{(16)}
\]

\inserted{\unsure{Prouvons|Prenons}} La \unsure{derniêre|dernière} relation (les deux premières étant évidentes).

Si on a\unsure{\deleted{vait}|vait|\deleted{ura}} $\Sigma_{12}=\varnothing$, comme $E_1,E_2\subset\Sigma_{12}$,
on aurait $E_1=\varnothing$, $E_2=\varnothing$, de sorte que $\unsure{\Pi|\pi}_1$ et $\Pi_2$ sont deux \uline{partitions} de $X$ en des parties non vides, i.e.
$\deleted{\Sigma_{\unsure{1|2}}^- = X\setminus\unsure{-|\Sigma}^{-}}$ on \unsure{aurait|avait}
\[
\deleted{\Sigma_1^+}\quad \Sigma_1^-=X\setminus\Sigma_1^+,\qquad \Sigma_2^-=X\setminus\Sigma_2^+
\]
\endgroup
